% =============================================================================
% Chapter 3 -- Design
% =============================================================================
\chapter{Design}
\label{ch:design}

This chapter justifies the methodological choices made before any
training run was launched. The project is dominated by experimental
design rather than software design --- the COMP390 handbook
explicitly permits collapsing this chapter and \Cref{ch:impl} into a
single ``Methodology'' chapter, but the two are kept separate here
to make the design rationale legible without the implementation
detail that would otherwise interleave with it.

\section{Overall Architecture}
\label{sec:arch}

\Cref{fig:arch} sketches the data flow and the module boundaries.
The COOM environment library wraps the ViZDoom binary and exposes a
sequence of Gymnasium-style \texttt{Env} objects. The Continual
Learning module composes those tasks into a sequence and drives a
single Soft Actor-Critic agent through them, with the regularisation
penalty injected by a method-specific subclass of a common base
class. Every training run writes a \texttt{progress.tsv} file and a
\texttt{config.json} snapshot of its hyperparameters; the
reproducibility layer downstream of training reads only those two
files, with no dependency on any in-process state.

\begin{figure}[t]
    \centering
    \fbox{\parbox{0.85\linewidth}{\centering\vspace{1.6cm}\textit{%
        ViZDoom $\rightarrow$ COOM environments
        $\rightarrow$ ContinualLearningEnv
        $\rightarrow$ SAC + regularisation method
        $\rightarrow$ progress.tsv / config.json
        $\rightarrow$ reproduction pipeline
        \vspace{1.6cm}}}}
    \caption{High-level system architecture. Each arrow is a
        narrowly typed interface; downstream modules cannot reach
        back across it.}
    \label{fig:arch}
\end{figure}

The deliberate design choice that this picture captures is the
hard separation between \emph{data generation} (the training
runs) and \emph{analysis} (everything that produces tables and
figures). The analysis pipeline is purely a function of the
recorded logs, so a finding can never depend on a state of the
training process that is not also written to disk. This makes the
study reproducible from logs alone, addresses
requirement~\ref{req:reproduce}, and was non-negotiable from the
start of the project.

\section{Experimental Design}
\label{sec:exp_design}

The headline study uses the CO4 task sequence
(\textsc{Chainsaw}~$\rightarrow$~\textsc{Raise the Roof}~$\rightarrow$%
~\textsc{Run and Gun}~$\rightarrow$~\textsc{Health Gathering}) at
$20{,}000$ training steps per task. Four design parameters define
the study; each is justified below.

\paragraph{Sequence choice.}
CO4 is the second half of the eight-task Cross-Objective sequence.
It is the smallest sequence available in COOM that still contains
heterogeneous objectives (combat, navigation, button-pressing,
collection), so any forgetting effect that depends on a change of
objective at a task boundary will be visible. CO8 was rejected
because the per-task training budget would have to be cut even
further to fit overnight on a CPU, taking the study out of the
regime where SAC is even partially trained. Cross-Domain (CD)
sequences were rejected because their tasks differ only in
visuals or dynamics, which makes forgetting a less stringent test of
the regularisation methods.

\paragraph{Per-task budget.}
The COOM paper uses $200{,}000$ steps per task. This dissertation
uses $20{,}000$ --- exactly one tenth. The factor was chosen so
that an entire CO4 run completes in under three hours on the
reference Apple-Silicon CPU, which is the largest budget that
permits five sequential runs to fit into a single overnight queue
(\Cref{nfr:queue}). The factor of ten is also large enough to be
unambiguously different from the published configuration: any
finding at this budget cannot be confounded with seed-level
variance in the published results. This tension --- between a
budget small enough to study and large enough to support
non-trivial learning --- is the single most important design
decision in the project. \Cref{ch:eval} reports that even at this
budget, three of the four methods do reach a non-trivial success
rate on at least one task, so the floor is correctly placed.

\paragraph{Method set.}
The first-stage method set comprises four methods: Fine-Tuning
(no regularisation), EWC, L2 anchoring, and MAS. The
unregularised control is essential --- a comparison among
regularised methods alone cannot answer whether regularisation
helps. EWC, L2 and MAS are the regularisation family reported
in the COOM paper. AGEM and VCL are the worst-performing
methods in the COOM ranking and would not change the ordering,
and so are dropped from the first-stage set. The hyperparameters
$\lambda_{\text{EWC}}=250$, $\lambda_{\text{L2}}=10^{5}$,
$\lambda_{\text{MAS}}=10^{4}$ are inherited verbatim from the
COOM paper. Re-tuning $\lambda$ at the new budget would have
collapsed the first-stage study into a sweep over a single
method and is left to second-stage work
(\Cref{sec:scope_update}).

\paragraph{Seed budget.}
EWC is run with three random seeds; the remaining first-stage
methods are single-seed. The seed allocation is set by the
clarity of the resulting figures rather than by a fixed
multi-seed protocol: three seeds are sufficient to draw a
visible band on the EWC curve, and the headline finding (the
inversion of the published ranking relative to Fine-Tuning) is
robust to the inter-seed spread observed in that band, which
removes the incentive to pay for additional seeds on the other
methods at the first stage. This is acknowledged as a partial
satisfaction of \ref{req:variance}; multi-seed runs for the
remaining methods are part of the second-stage scope below.

\section{Evaluation Metrics}
\label{sec:metrics}

The COOM paper reports three CL metrics: Average Performance,
Forgetting and Forward Transfer. Forgetting and Forward Transfer
both require evaluating the agent on past tasks during training,
which the COOM repository toggles via the \texttt{--test} flag at
some additional wall-clock cost. To preserve the overnight queue
budget the present study runs with \texttt{--no\_test}, and reports
\emph{training-time} success rate as a proxy. Two derived metrics
are then computed from the same logs:

\begin{itemize}
    \item \textbf{Last-five-epoch mean training success rate} per
        task, averaged across the CO4 tasks. This is the analogue
        of Average Performance; it is reported in \Cref{tab:avg_perf}.
    \item \textbf{Within-task drift}, defined as the peak
        training success during a task minus the last-five-epoch
        mean within the same task. This decouples the question
        ``does the agent end the task at the level it reached
        within the task?'' from the question ``does the agent
        recover when re-tested later?''. The first question can be
        answered from training logs alone; the second cannot, and
        is left to future work with \texttt{--test} enabled.
\end{itemize}

The substitution of training-time success for held-out evaluation
is a known threat to validity; it is discussed in
\Cref{sec:validity}.

\section{Alternatives Considered}
\label{sec:alts}

\begin{itemize}
    \item \emph{Longer per-task budget on cloud GPU.}\quad
        Rejected for cost reasons and to keep the first-stage
        study in the compute regime that the dissertation
        actually claims to characterise. Revisited in
        \Cref{sec:future} as the natural follow-on study.
    \item \emph{Replay-based methods (A-GEM, ClonEx-SAC) and
        architecture-based methods (PackNet).}\quad Deferred,
        not rejected. At the first stage the focus is on the
        regularisation family because mixing rehearsal or
        masking into the same comparison would make any
        observed effect a function of two design choices at
        once. After the first-stage results were reviewed with
        the supervisor, these higher-scoring methods were
        promoted into the second-stage method set
        (\Cref{sec:scope_update}); the rationale is that any
        improved EWC variant proposed by this dissertation
        should be benchmarked against the strongest published
        non-regularisation method (ClonEx-SAC), not only
        against its own regularisation peers.
    \item \emph{Smaller per-task budget ($5{,}000$ or
        $10{,}000$ steps).}\quad Considered and rejected at
        pilot time: at those budgets the agent failed to reach
        non-trivial success on any task, so the comparison
        would have been between four methods that all return
        zero. $20{,}000$ was the smallest budget at which the
        methods were still distinguishable.
    \item \emph{Held-out evaluation throughout
        training.}\quad Considered and rejected because enabling
        \texttt{--test} extends each run by roughly 20\,\%,
        which breaks the overnight queue budget for five
        sequential runs. Held-out Forgetting and Forward
        Transfer are part of the second-stage scope.
\end{itemize}

\section{Scope Update after Mid-Project Review}
\label{sec:scope_update}

After the first-stage results were reviewed with the supervisor
midway through the project, three changes to scope were agreed
on. The changes are recorded here, rather than absorbed silently
into earlier paragraphs, so that the marker can see the design
trajectory.

\paragraph{Sequence and budget unchanged.}
The CO4 sequence at $20{,}000$ steps per task was confirmed as
the right setting for both stages. Reducing the sequence
further or extending the per-task budget would have changed the
research question; keeping both fixed lets the second-stage
results be read directly against the first-stage table.

\paragraph{Seed budget set by figure clarity.}
The supervisor's guidance was that seeds should be allocated
by what the resulting figures need rather than by a uniform
protocol: enough seeds for the per-method band to be visible
and for the qualitative ranking to be defensible, but not so
many that the queue cost dominates the project. Three seeds for
each second-stage method is the working target; the figure
captions in \Cref{ch:eval} record the actual seed count for
each curve.

\paragraph{Method set broadened.}
The second-stage method set adds (a) cross-family baselines
--- ClonEx-SAC \citep{daniels2022clonex} as the strongest
published memory-based method on COOM and PackNet
(\citealp{mallya2018packnet}) as the strongest published
architecture-based method --- and (b) an \emph{improved EWC
variant} proposed by this dissertation, benchmarked
head-to-head against the original EWC of
\citet{kirkpatrick2017ewc} on identical hyperparameters
otherwise. The two additions address the supervisor's point
that ``the more baselines, the better'' for a credible
improved-method evaluation. Within the project window, the
improved EWC variant has been chosen, implemented and run
end-to-end (\Cref{sec:ewc_online_impl}, \Cref{sec:ewc_online}):
the variant is Online EWC \citep{schwarz2018progress}, which
replaces EWC's additive Fisher accumulation with an
exponential moving average. This choice was made on the basis
of the first-stage Fisher-on-noise diagnosis
(\Cref{sec:fishernoise}), which singles out the unbounded
growth of accumulated Fisher across tasks as the most direct
attack surface. The cross-family baselines and the multi-seed
extension of Online EWC are queued for follow-on work
(\Cref{sec:future}).
